
\section{平坦性的理想论判别法}

下面的定理将使我们能够对若干种环上的\trFlat{}模进行分类。它是\mbox{练习~23.1.C}（$M$ 是平坦 $A$-模当且仅当对一切 $A$-模 $N$ 都有 $\mathrm{Tor}^A_1(M,N)=0$）的精炼版本。

\textbf{24.3.1. \trhTheorem（平坦性的理想论判别法）。}\ ——\ \emph{$A$-\trModule{} $M$ 是\trFlat{}的，当且仅当对每个\trIdeal{} $I$ 都有 $\mathrm{Tor}^A_1(M,A/I)=0$。}

\textbf{24.3.2.}\ \emph{\trhRemark。}\quad 在给出证明之前，我们先作几条侧面的注记，它们或许有助于洞悉应当如何思考平坦性。\mbox{定理~24.3.1} 不借 $\mathrm{Tor}$ 理论之力来陈述也颇有收益。它等价于这样一个陈述：$M$ 是\trFlat{}的，当且仅当对一切\trIdeal{} $I\subset A$，$I\otimes_A M\to M$ 是\trInjectiveMap；你也可以以此面貌重新解读其证明。或许更好的说法是：$M$ 是\trFlat{}的，当且仅当对每个\trIdeal{} $I$，$I\otimes_A M\to IM$ 是\trIsomorphism。

平坦性常被非正式地描述成「连续变化的纤维」，而这可以按如下方式说得更加精确：$A$-\trModule{} $M$ 是\trFlat{}的，当且仅当它良好地限制到 $\Spec A$ 的诸\trClosedSubscheme{}上。更精确地说，在这个限制过程中我们所损失的东西——由「在 $Z$ 上消失」的元素组成的\trSubmodule{} $IM\subset M$——是容易理解的：它由诸元素 $i\otimes m$ 的形式线性组合构成，其间没有任何意外的关系——也就是说，它就是张量积 $I\otimes_A M$。这正是下面这道练习的内容，做它时你可以使用\mbox{定理~24.3.1}。

\textbf{24.3.A.}\ \hard\ \textbf{（练习：平坦性的方程判别法）}\quad 证明：$A$-\trModule{} $M$ 是\trFlat{}的，当且仅当对每个关系 $\sum_i a_i m_i=0$（其中 $a_i\in A$，$m_i\in M$），都存在 $m'_j\in M$ 与 $a_{ij}\in A$，使得对一切 $i$ 有 $\sum_j a_{ij}m'_j=m_i$，且对一切 $j$ 在 $A$ 中有 $\sum_i a_i a_{ij}=0$。（释读：每当 $M$ 的元素满足一条 $A$-线性关系时，这都是「因为」某些线性等式在 $A$ 中成立。）

\textbf{24.3.3.}\ \emph{不重要的注记。}\quad 在\mbox{定理~24.3.1} 的陈述中，只需对\trFinitelyGenerated{}的\trIdeal{}进行检验。这本质上是下面这条陈述的内容，你若愿意可以证一证：$A$-\trModule{} $M$ 是\trFlat{}的，当且仅当对一切\trFinitelyGenerated{}的\trIdeal{} $I$，自然映射 $I\otimes_A M\to M$ 是\trInjectiveMap。提示：若对某个非有限生成的\trIdeal{} $I$ 存在反例，就用它为某个\trFinitelyGenerated{}的\trIdeal{} $I'$ 找出另一个反例。

\textbf{24.3.4.}\ \emph{\mbox{定理~24.3.1}（平坦性的理想论判别法）的证明。}\quad 由\mbox{练习~23.1.C}，我们只需证明：对一切 $I$ 有 $\mathrm{Tor}^A_1(M,A/I)=0$ 蕴含对一切 $A$-\trModule{} $N$ 有 $\mathrm{Tor}^A_1(M,N)=0$，从而 $M$ 是\trFlat{}的。

我们先对 $N$ 的\trGenerator{} $a_1,\dots,a_n$ 的个数 $n$ 作归纳，证明对一切\trFinitelyGenerated{}的 $N$ 都有 $\mathrm{Tor}^A_1(M,N)=0$。基底情形（$n=1$，此时 $N\cong A/\mathrm{Ann}(a_1)$——好久以来头一次用到\trAnnihilatorIdeal）就是我们的假设。若 $n>1$，则 $Aa_n\cong A/\mathrm{Ann}(a_n)$ 是 $N$ 的一个\trSubmodule，而商 $Q$ 由 $a_1,\dots,a_{n-1}$ 的像生成，于是考虑
\[
0 \longrightarrow A/\mathrm{Ann}(a_n) \longrightarrow N \longrightarrow Q \longrightarrow 0
\]
的 $\mathrm{Tor}$ 长正合列中 $\mathrm{Tor}_1$ 的那一部分，结果便随之而来。

一般的 $N$ 的情形，我们用\trAbstractNonsense{}来处理。注意 $N$ 是它的诸\trFinitelyGenerated{}\trSubmodule{} $\{N_\alpha\}$ 的并。用时髦的语言说，这个并是一个\trFiltered\trColimit——任意两个有限生成\trSubmodule{}都含于某个有限生成\trSubmodule{}（具体地说，就是由它们生成的那个）。\trModule{}的\trFiltered\trColimit{}与上同调交换（\mbox{练习~1.5.M}），故 $\mathrm{Tor}_1(M,N)$ 就是诸 $\mathrm{Tor}_1(M,N_\alpha)=0$ 沿 $\alpha$ 取的\trColimit，于是证毕。\hfill$\square$

现在我们利用\mbox{定理~24.3.1}，对三种（类型的）环给出\trFlat{}模的显式刻画：\trPID、\trDualNumbers，以及某些局部环。

回想\mbox{观察~24.1.2}：平坦性蕴含\trTorsionFree。其逆对\trPID{}成立：

\textbf{24.3.B.（练习：\trPID{}上平坦＝无挠）}\quad 证明：\trPID{}上的模是\trFlat{}的当且仅当它是\trTorsionFree{}的。

\textbf{24.3.C.（练习：\trDVR{}上的平坦性）}\quad 设 $M$ 是\trDVR{} $A$ 上的模，$t$ 为其\trUniformizer。证明：$M$ 是\trFlat{}的当且仅当 $t$ 不是 $M$ 的\trZerodivisor，也就是说（沿用\mbox{练习~24.2.A} 定义的记号）$(M:t)=0$。（有限生成模的情形见\mbox{练习~24.1.B}。）这就给出了正则曲线上平坦性的一个简单而又十分重要的几何解读，我们将在 \S24.3.9 中讨论。

\textbf{24.3.D.（练习：\trDualNumbers{}上的平坦性）}\quad 证明：$M$ 在 $k[t]/(t^2)$ 上\trFlat{}当且仅当「乘以 $t$」映射 $M/tM\to tM$ 是\trIsomorphism。（这一事实在\trDeformationTheory{}和其他一些场合都很重要。）提示：$k[t]/(t^2)$ 只有三个\trIdeal。

\textbf{24.3.5. 重要的\trhTheorem（局部环上有限表现模：平坦＝自由＝投射）。}\ ——\ \emph{设 $(A,\mathfrak{m})$ 是\trLocalRing（不必\trNoether），$M$ 是\trFinitelyPresented{}的 $A$-\trModule。则以下诸条件等价。}
\begin{itemize}
\item[（a）]$M$ 自由。
\item[（b）]$M$ 是\trProjectiveObj{}的。
\item[（c）]$M$ 是\trFlat{}的。
\item[（d）]$\mathrm{Tor}^A_1(M,A/\mathfrak{m})=0$。
\end{itemize}

\textbf{24.3.6.}\ \emph{\trhRemark。}\quad \emph{警告}：局部环上的模可以平坦而不自由：$\mathbb{Q}$ 是\trFlat{}的 $\mathbb{Z}_{(p)}$-代数（$\mathbb{Z}_{(p)}$ 指 $\mathbb{Z}$ 在 $p$ 处的\trLocalization，而非\trPadics），因为一切局部化都是平坦的（\S24.1.1），但它不是自由的（你能看出为什么吗？）。

另外，非\trNoether{}阵营的人士或许乐于得知：只要稍费些功夫，「有限表现」可以减弱为「有限生成」：在下面的证明中出现有限表现的地方，改用 [Mat2, Thm.~7.10] 即可。

最后，\mbox{命题~24.4.5} 是\mbox{定理~24.3.5} 对分次环的一个十分近似的变体。

\emph{证明。}\quad 对任意环，自由模都是\trProjectiveObj{}的（\S23.2.1），而\trProjectiveObj{}模是\trFlat{}的（\mbox{练习~23.2.B}），并且\trFlat{}模 $M$ 满足对\emph{任意} $N$ 都有 $\mathrm{Tor}^A_1(M,N)=0$（\mbox{练习~23.1.C}），所以我们只需证明：若 $\mathrm{Tor}^A_1(M,A/\mathfrak{m})=0$，则 $M$ 自由。

（读到这里，你应该老远就嗅到\trNakayama{}要出场了。）此时 $M/\mathfrak{m}M$ 是域 $A/\mathfrak{m}$ 上的\trFiniteDimensional{}\trVectorSpace。选取 $M/\mathfrak{m}M$ 的一组基，并把它提升为 $M$ 中的元素 $m_1,\dots,m_n$。考虑由 $e_i\mapsto m_i$ 给出的映射 $A^{\oplus n}\to M$。我们将证明它是\trIsomorphism。由\trNakayama{}引理（见\mbox{练习~8.2.H}）它是满射：其像在模去极大理想后是全部 $M$，故其实就是全部。由于 $M$ 是有限表现的，据\mbox{引理~6.4.6}（「有限表现蕴含总是有限生成」），\trKernel{} $K$ 是有限生成的。把 $0\to K\to A^{\oplus n}\to M\to 0$ 与 $A/\mathfrak{m}$ 作张量。由于 $M$ \trFlat，所得结果仍然正合（\mbox{练习~24.2.C}）：
\[
0 \longrightarrow K/\mathfrak{m}K \longrightarrow (A/\mathfrak{m})^{\oplus n} \longrightarrow M/\mathfrak{m}M \longrightarrow 0.
\]
但 $(A/\mathfrak{m})^{\oplus n}\to M/\mathfrak{m}M$ 由构造即为\trIsomorphism，故 $K/\mathfrak{m}K=0$。而 $K$ 有限生成，据\trNakayama{}引理~8.2.9 得 $K=0$。\hfill$\square$

下面是一个立即可得且有用的推论——其实不过是一个几何解读罢了。

\textbf{24.3.7. \trhCorollary（有限表现层：平坦＝局部自由）。}\ ——\ \emph{$X$ 上的\trFinitelyPresented{}层 $\mathscr{F}$（在 $X$ 上）是\trFlat{}的，当且仅当它是\trLocallyFree{}的。}

\emph{证明。}\quad 有限表现层的局部自由性可以在茎上检验，见\mbox{练习~14.3.E}。\hfill$\square$

\textbf{24.3.E.（练习）}\quad 设 $\pi\colon X\to Y$ 是局部\trNoether{}概形之间的有限\trFlat{}\trMorphism，$\mathscr{F}$ 是 $X$ 上有限秩的\trLocallyFreeSheaf。
（a）证明 $\pi_\ast\mathscr{F}$ 是 $Y$ 上有限秩的\trLocallyFreeSheaf。
（b）若 $Y$ \trIrreducible{}且\trGenericPoint{}为 $\eta$，$\pi$ 在 $\eta$ 上方的\trDegree{}为 $n$，而 $\mathscr{F}$ 是秩 $r$ 的\trLocallyFreeSheaf，证明 $\pi_\ast\mathscr{F}$ 是秩 $nr$ 的\trLocallyFreeSheaf。（这可能比看上去更棘手。）

\textbf{24.3.8.}\ \emph{\trhRemark。}\quad 特别地，取 $\mathscr{F}=\OO_X$，便可见在局部\trNoether{}的设定下，\emph{有限\trFlat{}\trMorphism{}具有局部常值的\trDegree}。

\textbf{24.3.F.}\ \hard\ \textbf{（练习：\mbox{定理~24.3.5} 的有趣变体，但对我们并不重要）}\quad 设 $A$ 是环（不必局部），$M$ 是\trFinitelyPresented{}的 $A$-\trModule。证明：$M$ \trFlat{}当且仅当它是\trProjectiveObj{}的。提示：证明 $M$ 是\trProjectiveObj{}的当且仅当对每个极大理想 $\mathfrak{m}$，$M_\mathfrak{m}$ 自由。这一蕴含的较难方向要用到 $\mathrm{Hom}_{A_\mathfrak{m}}(M_\mathfrak{m},N_\mathfrak{m})=\mathrm{Hom}_A(M,N)_\mathfrak{m}$ 这一事实，它由\mbox{练习~1.5.H} 得出。（注记：存在非\trProjectiveObj{}的有限生成平坦模。它们必定不是有限表现的。例（不给出证明）：取 $A=\prod_{i=1}^{\infty}\mathbb{F}_2$，把它释读为函数 $\mathbb{Z}^{\ge 0}\to\mathbb{Z}/2$ 全体构成的环，并令 $M$ 为这些函数模去那些具有真支撑者——也就是说，在 $\mathbb{Z}^{\ge 0}$ 的几乎所有点处为零者——所得的模。）

\textbf{24.3.G.（练习）}\quad 证明下面这条有用的平坦性判别法：设 $\pi\colon X\to Y$ 是有限\trMorphism，$Y$ 既约且局部\trNoether。则 $\pi$ \trFlat{}当且仅当 $\pi_\ast\OO_X$ 是\trLocallyFree{}的，当且仅当 $\pi_\ast\OO_X$ 的\trRank{}局部常值（即 $\dim_{\kappa(\mathfrak{q})}(\pi_\ast\OO_X)_\mathfrak{q}\otimes\kappa(\mathfrak{q})$ 是 $\mathfrak{q}\in Y$ 的局部常值函数）。部分提示：\mbox{练习~14.3.K}。

\textbf{24.3.H.（练习）}\quad 证明结点的\trNormalization{}并不\trFlat（\mbox{图~8.4}）。提示：用\mbox{练习~24.3.G}。（这道练习还可以加强为：非平凡的\trNormalization\emph{绝不}\trFlat。）

\textbf{24.3.I.（练习）}\quad 在 $\mathbb{A}^4_k=\Spec k[w,x,y,z]$ 中，设 $X$ 是 $wx$-平面与 $yz$-平面的并：
\[
\tag{24.3.8.1}\label{eq:two-planes-union}
X=\Spec k[w,x,y,z]/(wy,wz,xy,xz).
\]
由 $k[a,b]\to k[w,x,y,z]$（$a\mapsto w-y$，$b\mapsto x-z$）给出的\trProjection{} $\mathbb{A}^4_k\to\mathbb{A}^2_k$ 限制为一个\trMorphism{} $X\to\mathbb{A}^2_k$。证明这个\trMorphism{}不是\trFlat{}的。

\textbf{24.3.9. 正则曲线上的平坦\trFamily。}\quad \mbox{练习~24.3.C} 给出了一个优雅的几何判别法，用以判别到正则\trCurve{}的\trMorphism{}何时\trFlat。

\textbf{24.3.J.（练习：正则曲线上平坦性的判别法）}\quad 设 $\pi\colon X\to Y$ 是从局部\trNoether{}概形到正则（局部\trNoether）\trCurve{}的\trMorphism。（对 $X$ 的局部\trNoether{}假设，是为了让我们能谈论它的\trAssociatedPoint。）证明：$\pi$ \trFlat{}当且仅当 $X$ 的一切\trAssociatedPoint{}都映到 $Y$ 的\trGenericPoint。（这是\mbox{练习~24.1.K}——平坦映射总把\trAssociatedPoint{}送到\trAssociatedPoint——的一个部分逆。与\mbox{练习~24.1.K} 一样，请尽管立刻把你的论证推广到 $X$ 上的一个\trCoherentSheaf{} $\mathscr{F}$。）

\textbf{24.3.10.}\quad 例如，从整（局部\trNoether）概形到正则\trCurve{}的非常值映射必定\trFlat。（另一个例子：到正规锥的\trDeformation——即带双星的 \S22.4.17 中讨论、\mbox{图~22.3} 所绘的那个——就是平坦的。）\mbox{练习~24.3.H}（及其后的评注）表明「正则」这一条件是必要的；\mbox{练习~24.3.I} 中两平面交于一点的例子则表明「维数 1」这一条件是必要的。

\textbf{24.3.11.}\ \hard\ \emph{不重要的注记：平坦性的一个赋值判别法。}\quad \mbox{练习~24.3.J} 表明，正则曲线上的平坦性在几何上是直观的（而且是「可以画出来」的）。它给出了一般情形的平坦性判别法：设 $\pi\colon X\to Y$ 是有限表现的\trMorphism。若 $\pi$ \trFlat，则对每个\trMorphism{} $Y'\to Y$——其中 $Y'$ 是某个\trDVR{}的 $\Spec$——都有 $\pi'\colon X\times_Y Y'\to Y'$ \trFlat，于是 $X\times_Y Y'$ 没有\trAssociatedPoint{}被映到 $Y'$ 的闭点。若 $Y$ 既约且局部\trNoether，则这也是充分条件；把它称作平坦性的一个\emph{\trValuativeCriteria}是十分合理的。（既约性是必要的：考虑\mbox{练习~24.1.G}(d)。）这给出了一种可视化平坦性的极好途径，你应当试着把它诉诸言语（或许要在下文学习了平坦极限之后）。（证明以及不设\trNoether{}假设的推广，见 [Gr-EGA, IV$_3$.11.8.1]。）

\textbf{24.3.12.}\ \emph{重访从\trCurve{}到正则\trCurve{}的射影\trMorphism{}的\trDegree。}\quad 正如\mbox{命题~16.3.5} 的陈述处所暗示的，我们现在能更好地理解：为什么从\trCurve{}到正则\trCurve{}的非常值射影\trMorphism{}具有良定的\trDegree，它可以通过取任意一点的\trInverseImage{}来确定（\S16.3.4）。（\mbox{例~10.3.4} 曾特别富有启发性。）缘由是：这样的映射由\mbox{练习~24.3.J} 是平坦的，而\trDegree{}由\mbox{注记~24.3.8} 是恒常的（另见\mbox{练习~24.3.G}）。此外，\mbox{练习~24.3.G} 还给出了\mbox{命题~16.3.5} 的一个新证明。

\textbf{24.3.13.}\ \hard\ \textbf{平坦极限。}\quad 这里是\mbox{练习~24.3.J} 的一个重要推论，我们可以非正式地把它陈述为：我们可以沿单参数\trFamily{}取平坦极限。更精确地说：设 $A$ 是\trDVR，$0$ 为 $\Spec A$ 的闭点，$\eta$ 为其\trGenericPoint。设 $X$ 是 $A$ 上的局部\trNoether{}概形，$Y$ 是 $X_\eta$ 的一个\trClosedSubscheme。令 $Y'$ 为 $Y$ 在 $X$ 中的\trSchemeTheoretic{}闭包。则由\mbox{练习~9.4.D}，$\pi\colon Y'\to\Spec A$ 在 $0$ 上方没有\trAssociatedPoint，故由\mbox{练习~24.3.J}，$Y'$ 在 $\Spec A$ 上平坦。类似地，设 $Z$ 是一维\trNoether{}概形，$0$ 是 $Z$ 的一个正则点，而 $\pi\colon X\to Z$ 是从局部\trNoether{}概形到 $Z$ 的\trMorphism。若 $Y$ 是 $\pi^{-1}(Z\setminus\{0\})$ 的一个\trClosedSubscheme，且在 $Z\setminus\{0\}$ 上平坦，而 $Y'$ 是 $Y$ 在 $X$ 中的\trSchemeTheoretic{}闭包，则 $Y'$ 在 $Z$ 上平坦。在这两种情形，闭包 $Y'|_0$ 常被称作 $Y$ 的\textbf{平坦极限}。（请尽管随意减弱 $X$ 上的\trNoether{}假设。）

\textbf{24.3.K.（练习）}\quad 假设（沿用上一段的语言）$A$ 是\trDVR，$X$ 是局部\trNoether{}的 $A$-概形，$Y$ 是一般纤维 $X_\eta$ 的一个\trClosedSubscheme。证明：$X$ 只有唯一的\trClosedSubscheme{} $Y'$，使得 $Y'|_\eta=Y$ 且 $Y'$ 在 $A$ 上平坦。

\textbf{24.3.L.（练习）}\quad 考虑 $\mathbb{A}^2_{\mathbb{C}}$ 中点对的\trFamily，其坐标由带坐标 $t$ 的 $\mathbb{A}^1_{\mathbb{C}}$ 参数化，给定为 $\{(0,0),(3t,4t)\}$。
（a）求 $t\to 0$ 时的平坦极限。（它将是一个长度为 $2$ 的\trSubscheme。这道练习的主旨是要你把其中的代数算到底。）
（b）求「$t\to\infty$」时的平坦极限。

\textbf{24.3.M.（练习）}\quad 设 $C=V(f(x,y))$ 是 $\mathbb{A}^2_k$ 中的一条光滑平面\trCurve，$p$ 是 $C$ 的一个\trClosedPoint，于是 $\OO_{C,p}$ 是\trDVR。考虑 $C$ 的点对的\trFamily，由 $\mathfrak{q}\in C\setminus\{p\}$ 参数化，给定为 $(p,\mathfrak{q})$。解释：为什么在 $p$ 上的平坦极限（即「$\mathfrak{q}\to p$ 时的极限」）可以与 $C$ 在 $p$ 处的切方向等同。释读：切线确实是割线的极限——而且无须 $\delta$ 与 $\epsilon$。

\textbf{24.3.N.（较难的练习：一个显式的平坦极限）}\quad 设域 $k$ 上有 $X=\mathbb{A}^3\times\mathbb{A}^1\to Z=\mathbb{A}^1$，其中 $\mathbb{A}^3$ 的坐标为 $x$、$y$、$z$，$\mathbb{A}^1$ 的坐标为 $t$。把 $Y$（在 $t=0$ 之外）定义为两条直线 $y=z=0$（即 $x$-轴）与 $x=z-t=0$（即平移了 $t$ 的 $y$-轴）的并。求 $t=0$ 处的平坦极限。（提示：（i）它并\emph{不是}两条坐标轴的并，尽管它包含这个并。平坦极限在结点处非既约，而且「\trFuzz」从它们所在的平面指了出来。（ii）$\langle y,z\rangle\langle x,z\rangle\ne\langle xy,z\rangle$。（iii）一旦你有了候选的平坦极限，务必检验它确实是平坦极限。（iv）若卡住了，请读下方的\mbox{例~24.3.14}。）

考虑上一个例子在射影情形下的版本：$\mathbb{P}^3$ 中两条直线退化而相交。极限由两条相交于一个结点的直线组成，结点处带有某种非既约结构。在两条直线靠拢之前，它们的整体截面空间是二维的。当它们靠拢时，它们的平坦极限是否也具有二维的整体截面空间，这一点并非一目了然。既约版本（相交于一点的两条直线之并）的整体截面空间是一维的，而非既约结构对整体截面空间的影响可能无法立刻看清。然而我们将看到，「上同调群在平坦极限中只会向上跳」，这是半连续性定理（\mbox{定理~25.1.1}）的一个推论（实际上也是它最主要的寓意）。

\textbf{24.3.14.}\ \extrahard\ \emph{平坦\trFamily{}中上同调群变化的一个例子。}\quad 我们可以用\mbox{练习~24.3.N} 的一个变体，来看一个上同调群真实跳动的例子。我们在代数闭域上工作，以免分心。在着手显式代数之前，先说总体想法。考虑 $\mathbb{P}^3$ 中的一条扭曲三次曲线 $C$。从随机的一点 $p\in\mathbb{P}^3$ 出发的\trProjection{} $pr_p$ 会把 $C$ 送到一条结点平面三次曲线。把这个\trProjection{}「动态地」图景化：选取坐标使得 $p$ 位于 $[1,0,0,0]$，并考虑映射 $\phi_t\colon[w,x,y,z]\mapsto[tw,x,y,z]$；$\phi_1$ 是 $\mathbb{P}^3$ 上的恒等，当 $t\ne 0$ 时 $\phi_t$ 是 $\mathbb{P}^3$ 的自同构，而 $\phi_0$ 正是那个\trProjection。当 $t\to 0$ 时，$\phi_t(C)$ 的极限将是一条结点三次曲线，其结点处的非既约结构，与\mbox{练习~24.3.N} 中两条直线靠拢时所见的「解析地相同」。（「解析地相同」这一说法，待我们在 \S28 中定义完备化之后便可使之精确。）

现在我们在实践中看一下。与其直接处理扭曲三次曲线，我们不妨用另一个见过类似图景的例子。考虑结点（仿射）平面三次曲线 $y^2=x^3+x^2$。它的\trNormalization（见\mbox{图~8.4}、\S8.3.5 之\mbox{例~(3)}、\mbox{练习~10.7.E} 等）是通过添加一个相应于 $y/x$ 的额外变量 $m$ 而得到的（这可以解读为胀开原点；见 \S22.4.3）。我们使用变量 $m$ 而非 $t$（\S8.3.5 中所用），为的是把 $t$ 留给平坦\trFamily{}的参数。

我们把结点\trCurve{} $C$ 画成位于 $\mathbb{A}^3=\Spec k[x,y,m]$ 的 $xy$-平面内，而其\trNormalization{} $\widetilde{C}$ \trProjection{}到它之上，$m=y/x$。$\widetilde{C}$ 的方程是什么？显然其中包括方程 $y^2=x^3+x^2$ 与 $y=mx$，但它们还不够——$m$-轴（即 $x=y=0$）也落在 $V(y^2-x^3-x^2,\,y-mx)$ 中。稍加思考（再加上本例中先前见过的代数）便会明白：我们还有第三个方程 $x=m^2-1$，它与 $y=mx$ 一起蕴含 $y^2=x^3+x^2$。现在方程终于够用了：$k[x,y,m]/(x-(m^2-1),\,y-mx)$ 是一个\trIntegralDomain，因为它显然\trIsomorphic{}于 $k[m]$。诚然，你应当认得
这正是\mbox{练习~10.7.E} 中出现的那套代数，于是我们把 $\widetilde{C}$ 描述为
\[
\tag{24.3.14.1}\label{eq:nodal-normalization}
\Spec k[x,y,m]/\bigl(x-(m^2-1),\,y-mx\bigr).
\]

接下来，我们想把「到 $\widetilde{C}\subset\mathbb{A}^3$ 之 $xy$-平面的动态投影」这一直觉形式化。我们把它图景化为如下：给定一点 $(x,y,z)$ 在时刻 $1$ 的位置，我们希望它在时刻 $t$ 位于 $(x,y,zt)$。在时刻 $t=1$，我们「从 $\widetilde{C}$ 出发」；而在时刻 $t=0$，我们（按集合论意义）得到 $C$。于是，在时刻 $t\ne 0$，曲线 $\widetilde{C}$ 被送到曲线
\[
\Spec k[x,y,m]/\bigl(x-((m/t)^2-1),\,y-(m/t)x\bigr)=\Spec k[x,y,m]/(m^2-t^2(x+1),\,ty-mx).
\]
从而 $\Spec k[t,t^{-1}]$ 上的\trFamily{}便是
\[
\Spec k[x,y,m,t,t^{-1}]/\bigl(m^2-t^2(x+1),\,ty-mx\bigr)\longrightarrow\Spec k[t,t^{-1}].
\]
我们之所以把 $t$ 反转，是因为到目前为止我们只处理非零的 $t$。对 $t\ne 0$，这当然是一个「好」的\trFamily，因而必定\trFlat。

\textbf{24.3.O.（练习）}\quad 请以尽可能无痛的方式验证这件事！提示：通过一次坐标变换，证明这个\trFamily{}「在 $\Spec k[t,t^{-1}]$ 上是常值的」，从而它是经由 $\Spec k[t,t^{-1}]\to\Spec k$、从
\[
\Spec k[x,y,M]/\bigl(M^2-(x+1),\,y-Mx\bigr)\longrightarrow\Spec k
\]
拉回（以某种有待你自己想明白的方式）而得到的；而后者由\mbox{陷阱问题~24.1.A}(a) 是\trFlat{}的。

现在我们在 $\Spec k[x,y,m,t]\to\mathbb{A}^1=\Spec k[t]$ 中把这个\trFamily{}在 $t=0$ 上的平坦极限算出来。我们首先希望这个平坦\trFamily{}由我们已经写下的方程给出：
\[
\Spec k[x,y,m,t]/\bigl(m^2-t^2(x+1),\,ty-mx\bigr).
\]
但它\emph{不}在 $\mathbb{A}^1=\Spec k[t]$ 上平坦，因为纤维维数会跳动（\S24.5.5）：把 $t=0$ 代入方程（即取 $0\in\mathbb{A}^1$ 上的纤维），我们得到 $\Spec k[x,y,m]/(m^2,mx)$。按集合论意义，这是整个 $xy$-平面（$m=0$），其维数当然是 $2$。为便于后文引用，请注意这个「假极限」作为概形就是 $xy$-平面，\emph{带有些许沿 $y$-轴的非既约结构}。（这或许会让你想起\mbox{图~4.4}。）

所以我们至少还漏了一个方程。至于漏的是哪个方程，有一条线索：方程 $y^2=x^3+x^2$ 在 $t\ne 0$ 时显然成立，但对于我们那个朴素的极限方案 $m^2=mx=0$ 却\emph{不}成立。事后来看，我们在写 \eqref{eq:nodal-normalization} 时本不该把这条原始方程丢弃。于是把它放回去，我们便有了描述 $\mathbb{A}^1$ 上平坦\trFamily{}的第二个希望：
\[
\Spec k[x,y,m,t]/\bigl(m^2-t^2(x+1),\,ty-mx,\,y^2-x^3-x^2\bigr)\longrightarrow\Spec k[t].
\]
当 $t=0$ 时，我们得到 $\Spec k[x,y,m]/(m^2,mx,\,y^2-x^3-x^2)$，于是它通过了「无维数跳动」的平坦性判别法（又是 \S24.5.5），而这只是必要条件、并非充分。然而遗憾的是，它并不\trFlat，而且这一点\footnote{原书此处作「and this not so easy to see」，当脱「is」一词，径补。——译者注}并非那么容易看出。

\textbf{24.3.P.（练习）}\quad 证明：在环
\[
k[x,y,m,t]/\bigl(m^2-t^2(x+1),\,ty-mx,\,y^2-x^3-x^2\bigr)
\]
中，$my-tx(x+1)\ne 0$，但 $t\bigl(my-tx(x+1)\bigr)=0$。

这提示了我们应当添加哪个方程：$my-tx(x+1)=0$。为方便计，令
\[
A=k[x,y,m,t]/\bigl(ty-mx,\,m^2-t^2(x+1),\,y^2-x^2(x+1),\,my-tx(x+1)\bigr).
\]
我们接下来希望 $\Spec A\to\Spec k[t]$ 在 $t=0$ 处\trFlat。我们可以用\mbox{练习~24.3.B} 来证明这一点，办法是证明 $t$ 不是 $A$ 的\trZerodivisor。而我们做到这件事的途径，是给出 $A$ 中元素的一个「\trNormalForm」。

\textbf{24.3.Q.（练习）}\quad 证明：$A$ 的每个元素都可以唯一地写成 $f(t)m+g(x,t)y+h(x,t)$ 的形式，其中 $f$、$g$、$h$ 是多项式。然后证明 $t$ 不是 $A$ 的\trZerodivisor，并由此得出 $\Spec A\to\mathbb{A}^1$ 确实\trFlat。

特别地，$t=0$ 时的平坦极限由下式给出：
\[
\Spec A/(t)=\Spec k[x,y,m]/(m^2,\,mx,\,my,\,y^2-x^2-x^3).
\]

\textbf{24.3.R.（练习）}\quad 证明：这个平坦极限非既约，而且「其非既约性的\trLength{}为 $1$，支撑在原点处」。更精确地说，若 $X=\Spec A/(t)$，证明 $\mathscr{I}_{X^{\mathrm{red}}}$ 是一个\trSkyscraperSheaf，取值为 $k$，支撑在原点。请为这个平坦极限 $X$ 画一张草图。

\textbf{24.3.15.}\quad 注意：$X$ 上有一个非零的整体函数，由 $m$ 给出，其\trSupport{}在原点（也就是说，在原点之外为 $0$）。

现在我们利用这个例子来构造一个行为有趣的射影例子。取这个例子的\trProjectiveCompletion，便得到 $\mathbb{P}^3$ 中诸三次曲线的一个\trFamily，它退化为一条带非既约点的结点三次曲线 $C$。

\textbf{24.3.S.（练习）}\quad 把这件事做了：精确地描述这个\trFamily（位于 $\mathbb{P}^3\times\mathbb{A}^1$ 中）。

取相应于
\[
0\longrightarrow\mathscr{I}_{C^{\mathrm{red}}}\longrightarrow\OO_C\longrightarrow\OO_{C^{\mathrm{red}}}\longrightarrow 0
\]
的\trLongExactSequence，便得到
\begin{equation*}
\begin{tikzcd}[row sep=2.6em, column sep=2.2em]
0 \ar[r] & \mathrm{H}^0(C,\mathscr{I}_{C^{\mathrm{red}}}) \ar[r,"\alpha"] & \mathrm{H}^0(C,\OO_C) \ar[r] & \mathrm{H}^0(C,\OO_{C^{\mathrm{red}}}) & \\
 \ar[r] & \mathrm{H}^1(C,\mathscr{I}_{C^{\mathrm{red}}}) \ar[r] & \mathrm{H}^1(C,\OO_C) \ar[r] & \mathrm{H}^1(C,\OO_{C^{\mathrm{red}}}) & \\
 \ar[r] & \mathrm{H}^2(C,\mathscr{I}_{C^{\mathrm{red}}}). & & &
\end{tikzcd}
\end{equation*}
我们有 $\mathrm{H}^1(C,\mathscr{I}_{C^{\mathrm{red}}})=\mathrm{H}^2(C,\mathscr{I}_{C^{\mathrm{red}}})=0$，因为 $\mathscr{I}_{C^{\mathrm{red}}}$ 的\trSupport{}在维数 $0$ 处（由维数\trCohomology{}消失，\mbox{定理~18.2.6}）。另外，$\mathrm{H}^1(C^{\mathrm{red}},\OO_{C^{\mathrm{red}}})=\mathrm{H}^1(C,\OO_{C^{\mathrm{red}}})$（上同调的性质 \textbf{(v)}；见 \S18.1）。（既约的）结点三次曲线 $C^{\mathrm{red}}$ 满足 $h^0(\OO)=1$（\S11.4.7）与 $h^1(\OO)=1$（三次平面曲线的\trGenus{}为 $1$，\eqref{eq:genus-plane-curve}）。还有，$h^0(C,\mathscr{I}_{C^{\mathrm{red}}})=1$，即上文所观察者。最后，$\alpha$ 不是 $0$，因为 $C$ 上存在一个在 $C^{\mathrm{red}}$ 上消失的非零函数（\S24.3.15——请说服自己：这个函数可以从仿射片 $\Spec A$ 延拓到\trProjectiveCompletion）。

利用这条\trLongExactSequence，我们得出 $h^0(C,\OO_C)=2$ 与 $h^1(C,\OO_C)=1$。于是在这个例子中我们看到：\trFamily{}的一般成员（扭曲三次曲线\trIsomorphic{}于 $\mathbb{P}^1$）满足 $(h^0(\OO),h^1(\OO))=(1,0)$，而特殊成员（
平坦极限）则满足 $(h^0(\OO),h^1(\OO))=(2,1)$。注意：两个\trCohomology{}群都跳动了，然而\trEulerCharacteristic{}却保持不变。第一种行为，正如\mbox{练习~24.3.N} 之后所述，是半连续性定理（\mbox{定理~25.1.1}）的一个例子；第二种行为——平坦\trFamily{}中\trEulerCharacteristic{}的恒常性——则是 \S24.7 的主题。（这个例子的极限是奇异的，这绝非巧合；见 \S25.1.3。）
